\chapter{Concepts Fondamentaux}
\begin{spacing}{1.2}
\minitoc
\thispagestyle{MyStyle}
\end{spacing}
\clearpage

\section{Introduction}

Ce chapitre a pour objectif de présenter les concepts techniques qui constituent la
base de compréhension du projet. Avant d'aborder l'architecture et l'implémentation
de la plateforme, il est nécessaire de clarifier les notions fondamentales sur
lesquelles elle repose~: le DevSecOps, le Cloud Computing, l'Infrastructure as Code,
les pipelines CI/CD, la conteneurisation, et l'intégration de l'intelligence
artificielle dans les environnements d'automatisation. Ces concepts ne sont pas
simplement des buzzwords~; ils définissent concrètement la manière dont les systèmes
logiciels modernes sont construits et exploités.

\section{DevSecOps : culture et pratiques}

\subsection{Origines et définition}

Le terme DevSecOps est issu de la fusion de trois disciplines~: le développement
logiciel (\textit{Dev}), les opérations système (\textit{Ops}), et la sécurité
(\textit{Sec}). Il représente une évolution naturelle du mouvement DevOps apparu dans
les années 2010, dont l'objectif était de rapprocher les équipes de développement et
d'exploitation pour accélérer les cycles de livraison. L'ajout de la sécurité comme
composante native~--- plutôt que comme une étape finale de validation~--- est ce qui
définit le DevSecOps.

Concrètement, l'idée est d'intégrer les contrôles de sécurité directement dans le
pipeline de développement et de déploiement, à chaque étape, plutôt que de traiter
la sécurité comme une phase distincte. On parle de \textit{shift-left}~: les
vérifications de vulnérabilités, les analyses de code statique et les tests de
conformité sont effectués le plus tôt possible dans le cycle de vie logiciel.

\subsection{Principes clés}

Plusieurs principes structurent les pratiques DevSecOps. Le premier est
l'\textbf{automatisation}~: les contrôles manuels ne peuvent pas suivre la cadence
des déploiements continus. Il faut donc automatiser les tests, les scans de
vulnérabilités, et les validations de configuration. Le deuxième principe est la
\textbf{traçabilité}~: chaque artefact déployé doit pouvoir être retracé jusqu'à son
code source, ses dépendances et ses tests. Enfin, la \textbf{responsabilité partagée}
implique que la sécurité n'est pas le problème d'une seule équipe~--- chaque
développeur, opérateur et architecte est concerné.

\subsection{DevSecOps et intelligence artificielle}

L'émergence des outils d'IA appliqués au DevSecOps a ouvert de nouvelles
perspectives. Les modèles de langage peuvent analyser des logs d'erreur, suggérer
des corrections de configuration, ou détecter des patterns anormaux dans les
pipelines. C'est précisément dans cette direction que s'inscrit notre projet~: tirer
parti des capacités des LLM pour automatiser des tâches qui requièrent habituellement
une expertise humaine avancée.

\section{Cloud Computing et Infrastructure as Code}

\subsection{Cloud Computing : rappel des modèles}

Le Cloud Computing désigne la mise à disposition de ressources informatiques~---
serveurs, stockage, bases de données, réseau~--- via internet, à la demande et
généralement à la consommation. On distingue trois grands modèles de service~:
l'IaaS (\textit{Infrastructure as a Service}), le PaaS (\textit{Platform as a
Service}) et le SaaS (\textit{Software as a Service}). Dans le contexte de ce
projet, c'est principalement l'IaaS qui est visé, avec AWS comme fournisseur de
référence.

Les architectures cloud modernes sont dites \textit{cloud-native} lorsqu'elles sont
conçues pour exploiter les caractéristiques propres au cloud~: élasticité, haute
disponibilité, distribution géographique, et facturation à l'usage. Ce modèle impose
des choix techniques spécifiques, notamment l'utilisation de conteneurs, de
microservices et d'orchestrateurs.

\subsection{Infrastructure as Code}

L'Infrastructure as Code (IaC) consiste à décrire l'infrastructure d'un système sous
forme de fichiers de configuration versionnables, lisibles et automatisables.
Plutôt que de configurer des serveurs manuellement via une interface graphique, on
écrit du code qui décrit l'état désiré de l'infrastructure, et un outil se charge
de l'appliquer.

\textbf{Terraform}, développé par HashiCorp, est l'outil IaC le plus répandu dans
les environnements cloud. Il utilise un langage déclaratif, le HCL
(\textit{HashiCorp Configuration Language}), pour décrire des ressources cloud AWS,
Azure, GCP, et bien d'autres. Son approche déclarative signifie que l'on décrit
\textit{ce qu'on veut}, non \textit{comment y arriver}~--- Terraform se charge du
plan d'exécution.

Un point important est la gestion de l'\textbf{état}~: Terraform maintient un
fichier d'état (\texttt{terraform.tfstate}) qui représente la correspondance entre
ce qui est décrit dans le code et ce qui existe réellement dans le cloud. Cette
notion est fondamentale pour comprendre pourquoi la génération automatique de fichiers
Terraform doit être rigoureuse et cohérente.

\subsection{Estimation des coûts cloud}

L'un des défis majeurs des infrastructures cloud est la maîtrise des coûts. Selon la
nature des ressources provisionnées~--- instances EC2, bases de données RDS, services
S3, load balancers~--- la facture mensuelle peut varier considérablement. Les
pratiques FinOps, qui combinent ingénierie financière et exploitation cloud,
recommandent d'anticiper ces coûts dès la phase de conception. C'est pourquoi la
plateforme développée dans ce projet intègre un module d'estimation des coûts AWS,
permettant de comparer différentes stratégies d'architecture avant tout déploiement.

\section{Intégration et Déploiement Continus (CI/CD)}

\subsection{Principes du CI/CD}

L'intégration continue (CI) désigne la pratique qui consiste à intégrer fréquemment
les modifications de code dans un dépôt partagé, en déclenchant automatiquement une
série de vérifications~: compilation, tests unitaires, analyse statique. L'objectif
est de détecter les régressions le plus tôt possible, avant que le code ne soit
fusionné dans la branche principale.

Le déploiement continu (CD) étend ce principe à la livraison et au déploiement.
Une fois le code validé, il peut être automatiquement déployé en environnement de
préproduction, voire de production, sans intervention manuelle. Cette automatisation
réduit les délais de mise en production et limite les erreurs humaines liées aux
déploiements manuels.

\subsection{Structure d'un pipeline CI/CD}

Un pipeline CI/CD est typiquement composé de plusieurs étapes enchaînées~:

\begin{itemize}
    \item \textbf{Source~:} déclenchement du pipeline à chaque \textit{push} ou
    \textit{pull request} sur le dépôt.
    \item \textbf{Build~:} compilation du code et construction des images Docker si
    applicable.
    \item \textbf{Test~:} exécution des tests unitaires, d'intégration, et des
    analyses de sécurité (\textit{SAST}, \textit{DAST}).
    \item \textbf{Package~:} création des artefacts déployables (images, archives).
    \item \textbf{Deploy~:} déploiement automatisé en environnement cible via des
    outils comme Kubernetes ou Terraform.
\end{itemize}

Les outils les plus utilisés pour implémenter ces pipelines sont GitHub Actions,
GitLab CI, Jenkins et ArgoCD pour les déploiements continus sur Kubernetes.

\subsection{CI/CD et DevSecOps}

Dans une approche DevSecOps, les pipelines CI/CD ne se contentent pas de tester le
code~; ils intègrent des étapes de vérification de sécurité automatisée. Cela inclut
l'analyse des dépendances pour détecter des vulnérabilités connues, le scan des
images Docker pour identifier des failles dans les couches de base, et la validation
des configurations d'infrastructure avant déploiement. Ces vérifications, si elles
sont bien intégrées, ne ralentissent pas le pipeline~--- elles le sécurisent.

\section{Conteneurisation et Orchestration}

\subsection{Docker et la conteneurisation}

Un conteneur est une unité d'exécution légère et portable qui encapsule une
application avec toutes ses dépendances~: bibliothèques, configuration, environnement
d'exécution. Contrairement à une machine virtuelle, un conteneur partage le noyau du
système hôte, ce qui le rend beaucoup plus léger et plus rapide à démarrer.

\textbf{Docker} est la technologie de conteneurisation qui a largement popularisé ce
paradigme. Un \texttt{Dockerfile} décrit, sous forme d'instructions séquentielles,
comment construire l'image d'un conteneur. Cette image est ensuite instanciée en
conteneur sur n'importe quelle machine disposant du moteur Docker, garantissant la
reproductibilité de l'environnement.

La qualité d'un Dockerfile a des implications directes sur la sécurité, la
performance et la maintenabilité d'une application~: taille de l'image, choix de
l'image de base, droits d'exécution, gestion des variables d'environnement. C'est
précisément pourquoi la génération automatique de Dockerfiles par un système IA
nécessite des contraintes strictes et une validation systématique.

\subsection{Kubernetes et l'orchestration de conteneurs}

Lorsque le nombre de conteneurs à gérer augmente, la gestion manuelle devient
impraticable. \textbf{Kubernetes} est la plateforme d'orchestration de conteneurs qui
s'est imposée comme le standard de l'industrie pour gérer des flottes de conteneurs à
grande échelle.

Kubernetes prend en charge le déploiement, la mise à l'échelle automatique, la
répartition de charge, la récupération sur panne, et la gestion des mises à jour sans
interruption de service. Il introduit une série d'abstractions~: les
\textit{Pods} (groupes de conteneurs), les \textit{Services} (exposition réseau), les
\textit{Deployments} (gestion du cycle de vie), et les \textit{ConfigMaps} et
\textit{Secrets} (gestion de la configuration). Ces ressources sont déclarées sous
forme de fichiers YAML, qui s'inscrivent naturellement dans une approche IaC.

\subsection{Architecture cloud-native}

L'association de Docker, Kubernetes et Terraform constitue la base d'une architecture
\textit{cloud-native} moderne. Docker assure la portabilité des applications,
Kubernetes leur orchestration, et Terraform le provisionnement de l'infrastructure
sous-jacente. Cette combinaison permet de gérer des systèmes distribués de manière
reproductible, automatisée et versionnée.

\section{Intelligence Artificielle appliquée à l'automatisation DevSecOps}

\subsection{Large Language Models et génération de code}

Les \textit{Large Language Models} (LLM) sont des modèles de langage entraînés sur
de très grands corpus de textes, capables de comprendre et de générer du texte dans
un très large spectre de contextes. Des modèles comme GPT d'OpenAI ont démontré des
capacités remarquables dans des tâches de génération de code, d'analyse de logs,
de suggestion de correctifs et de raisonnement technique.

Dans le contexte DevSecOps, les LLM peuvent être utilisés pour générer
automatiquement des fichiers de configuration (Dockerfile, YAML Kubernetes, HCL
Terraform), analyser des messages d'erreur et proposer des corrections, ou encore
comparer des architectures cloud selon différents critères. Ces capacités sont
cependant conditionnées par la qualité des prompts et du contexte fourni au modèle.

\subsection{Architectures multi-agents}

Un système multi-agents est composé de plusieurs agents autonomes, chacun spécialisé
dans une tâche précise, et capables de collaborer pour accomplir un objectif global.
Cette approche est particulièrement adaptée aux tâches complexes qui ne peuvent pas
être résolues efficacement par un seul modèle générique.

Dans le cadre de l'automatisation DevSecOps, une architecture multi-agents permet de
décomposer le workflow en étapes spécialisées~: un agent analyse le code source, un
autre prend des décisions architecturales, un troisième génère les fichiers Terraform,
un quatrième valide les résultats et un cinquième détecte et corrige les erreurs. La
spécialisation de chaque agent améliore la qualité des sorties et facilite la
maintenance du système.

\subsection{Retrieval-Augmented Generation (RAG)}

Le RAG (\textit{Retrieval-Augmented Generation}) est une technique qui enrichit les
requêtes envoyées à un LLM en y adjoignant des documents pertinents extraits d'une
base de connaissances. Plutôt que de se fier uniquement aux connaissances encodées
dans les paramètres du modèle, le système récupère dynamiquement des informations
contextuelles et les intègre dans le prompt.

Cette approche est particulièrement utile dans un contexte DevSecOps, où les bonnes
pratiques évoluent rapidement et où des connaissances spécifiques~--- bonnes pratiques
Terraform, patterns de sécurité Docker, configurations Kubernetes recommandées~---
doivent être intégrées dans les générations. Un système RAG spécialisé par domaine
technique permet de produire des artefacts mieux alignés sur les standards actuels,
avec moins d'hallucinations.

\subsection{Structured Constrained Generation}

La génération libre par LLM, bien que flexible, souffre d'un problème fondamental~:
l'absence de structure garantie. Pour un Dockerfile ou un fichier Terraform, une
erreur de syntaxe ou une incohérence entre sections peut rendre l'artefact
inutilisable. La \textit{Structured Constrained Generation} est une approche qui
découpe la génération en sections indépendantes, chacune guidée par un
\textit{blueprint} structuré qui contraint la forme et le contenu de la sortie. Les
sections sont ensuite assemblées de manière déterministe. Cette technique, adoptée
dans notre projet, améliore significativement la fiabilité et la validité des
artefacts générés.

\section{Sécurité dans les environnements DevOps}

\subsection{Sécurité des images et des conteneurs}

L'une des surfaces d'attaque les plus négligées dans les environnements conteneurisés
est l'image Docker elle-même. Une image construite sur une base obsolète ou contenant
des bibliothèques vulnérables expose l'application à des risques réels. Les bonnes
pratiques recommandent d'utiliser des images de base minimales (\texttt{alpine},
\texttt{distroless}), de ne pas exécuter les processus en tant que \texttt{root}, et
de scanner les images avec des outils comme Trivy ou Snyk avant déploiement.

\subsection{Sécurité de l'Infrastructure as Code}

Les fichiers IaC peuvent eux-mêmes contenir des erreurs de sécurité~: des ports
ouverts inutilement, des buckets S3 publics, des permissions IAM trop larges. Des
outils comme \textit{tfsec} ou \textit{Checkov} permettent d'analyser
statiquement les fichiers Terraform et d'identifier ces problèmes avant qu'ils ne
soient appliqués en production. Intégrer ces vérifications dans le pipeline CI/CD,
avant le \texttt{terraform apply}, est une pratique DevSecOps fondamentale.

\subsection{Gestion des secrets}

Dans un pipeline d'automatisation, la gestion des credentials est un enjeu critique.
Les mots de passe, clés API, et tokens d'accès ne doivent jamais apparaître en clair
dans les fichiers de configuration ou les dépôts de code. Des solutions comme
AWS Secrets Manager, HashiCorp Vault ou les \textit{Secrets} Kubernetes permettent
de centraliser et de sécuriser ces informations, en les injectant dynamiquement dans
les environnements d'exécution.

\section{Conclusion}

Ce chapitre a posé les bases conceptuelles nécessaires à la compréhension du projet.
Le DevSecOps, l'Infrastructure as Code, les pipelines CI/CD, la conteneurisation
avec Docker et Kubernetes, et l'utilisation des LLM dans des architectures
multi-agents forment le socle technologique sur lequel repose la plateforme
développée. Les notions de RAG et de \textit{Structured Constrained Generation}
expliquent les choix d'architecture IA adoptés pour garantir la fiabilité des
artefacts produits. Le chapitre suivant sera consacré à l'analyse des besoins et à
la conception détaillée du système.
